% Part: sets-functions-relations
% Chapter: relations
% Section: relations-as-sets

\documentclass[../../../include/open-logic-section]{subfiles}

\begin{document}

\olfileid{sfr}{rel}{set}
\olsection{关系作为集合}

\begin{explain}
在 \olref[sfr][set][imp]{sec} 中，我们提到了一些重要的集合：
$\Nat$、$\Int$、$\Rat$、$\Real$。你肯定记得这些集合元素之间的一些有趣关系。
例如，每个这样的集合上都有一个完全标准的\emph{序关系}。
此外，还有每个对象对自身且仅对自身具有的\emph{恒等}关系。
我们还会遇到更多有趣的关系，可能的关系甚至更多。
但在回顾它们之前，我们先要指出，我们可以把关系看作一种特殊的集合。

为此，请回忆 \olref[sfr][set][pai]{sec} 中的两点。
首先，回忆\emph{有序对}的概念：给定 $a$ 和 $b$，我们可以构成~$\tuple{a, b}$。
重要的是，这里元素的顺序\emph{确实}重要。所以如果 $a \neq b$，那么 $\tuple{a, b} \neq \tuple{b, a}$。
（这与无序对（即二元集合）形成对比，在无序对中 $\{a, b\}=\{b, a\}$。）
其次，回忆\emph{笛卡尔积}的概念：如果 $A$ 和 $B$ 是集合，那么我们可以构成~$A \times B$，
即所有满足 $x \in A$ 且 $y \in B$ 的序对 $\tuple{x, y}$ 的集合。
特别地，$A^{2}= A \times A$ 是来自~$A$ 的所有有序对的集合。

现在我们来考虑集合上的一个特定关系：自然数集合~$\Nat$ 上的 $<$-关系。
考虑所有满足 $n<m$ 的数对 $\tuple{n, m}$ 的集合，即
\[
  R=\Setabs{\tuple{n, m}}{n, m \in \Nat \text{ 且 } n<m}.
\]
$n$ 小于 $m$ 与序对 $\tuple{n, m}$ 属于 $R$ 之间有着密切的联系，即：
\[
      n<m\text{ 当且仅当 }\tuple{n, m} \in R.
\]
事实上，我们可以在不丢失任何信息的情况下，将集合 $R$ 视为 $\Nat$ 上的 $<$-关系\emph{本身}。

同样，我们可以为数之间的任意关系构造一个 $\Nat^{2}$ 的子集。
反过来，给定任意数对集合 $S \subseteq \Nat^{2}$，都有一个对应的数之间的关系，
即 $n$ 对 $m$ 具有该关系当且仅当 $\tuple{n, m} \in S$。这就为下面的定义提供了依据：
\end{explain}

\begin{defn}[二元关系]
集合 $A$ 上的一个\emph{二元关系}是~$A^{2}$ 的一个子集。
如果 $R \subseteq A^{2}$ 是~$A$ 上的一个二元关系，且 $x, y \in A$，
我们有时将 $\tuple{x, y} \in R$ 记作 $Rxy$（或 $xRy$）。
\end{defn}

\begin{ex}
  \ollabel{relations}
自然数对构成的集合 $\Nat^{2}$ 可以排列成如下的二维矩阵：
\[
  \begin{array}{ccccc}
  \mathbf{\tuple{ 0,0 }} & \tuple{ 0,1 } &
    \tuple{ 0,2 } & \tuple{ 0,3 } & \ldots\\
  \tuple{ 1,0 } & \mathbf{\tuple{ 1,1 }} &
    \tuple{ 1,2 } & \tuple{ 1,3 } & \ldots\\
  \tuple{ 2,0 } & \tuple{ 2,1 } &
    \mathbf{\tuple{ 2,2 }} & \tuple{ 2,3 } & \ldots\\
  \tuple{ 3,0 } & \tuple{ 3,1 } & \tuple{ 3,2 } &
    \mathbf{\tuple{ 3,3 }} & \ldots\\
  \vdots & \vdots & \vdots & \vdots & \mathbf{\ddots}
  \end{array}
\]

这里我们把对角线加粗显示，因为 $\Nat^2$ 中位于对角线上的那些序对构成的子集，即
\[
  \{\tuple{0,0 }, \tuple{ 1,1 }, \tuple{ 2,2 }, \dots\},
  \]
就是 $\Nat$ 上的\emph{恒等关系}。
（由于恒等关系很常用，让我们定义 $\Id{A}=\Setabs{\tuple{ x,x }}{x \in
A}$ 为任何集合 $A$ 上的恒等关系。）
所有位于对角线上方的序对构成的子集，即
\[
  L = \{\tuple{ 0,1 },\tuple{ 0,2 },\ldots,\tuple{ 1,2 },
  \tuple{ 1,3 }, \dots, \tuple{ 2,3 }, \tuple{ 2,4 },\ldots\},
\]
是\emph{小于}关系，即 $Lnm$ 当且仅当 $n<m$。所有位于对角线下方的序对构成的子集，即
\[
  G=\{ \tuple{ 1,0 },\tuple{ 2,0 },\tuple{
    2,1 }, \tuple{ 3,0 },\tuple{ 3,1 },\tuple{ 3,2 }, \dots\},
\]
是\emph{大于}关系，即 $Gnm$ 当且仅当 $n>m$。$L$ 与 $I$ 的并集，我们可以称之为 $K=L\cup I$，是\emph{小于等于}关系：
$Knm$ 当且仅当 $n \le m$。类似地，$H=G \cup I$ 是\emph{大于等于关系}。
这些关系 $L$、$G$、$K$ 和 $H$ 是一类特殊的关系，称为\emph{序}。
$L$ 和 $G$ 具有这样的性质：没有任何数对自身具有 $L$ 或 $G$ 关系（即，对所有 $n$，既非 $Lnn$ 也非 $Gnn$）。
具有这种性质的关系被称为\emph{反自反的}，并且，如果它们同时也是序，则被称为\emph{严格序}。
\end{ex}

\begin{explain}
尽管序和恒等是重要且自然的关系，但应该强调的是，根据我们的定义，\emph{任何} $A^{2}$ 的子集都是~$A$ 上的一个关系，
无论它看起来多么不自然或多么刻意。
特别地，$\emptyset$ 是任何集合上的一个关系（\emph{空关系}，没有任何元素对具有该关系），
而 $A^{2}$~本身也是~$A$ 上的一个关系，称为\emph{全关系}（所有元素对都具有该关系）。
但即使是像
$E=\Setabs{\tuple{n, m}}{n>5 \text{ or } m \times n \ge 34}$
这样的集合也算作一个关系。
\end{explain}

\begin{prob}
  列出集合
  $\Pow{\{a, b, c\}}$
  上关系 $\subseteq$ 的所有元素。
\end{prob}

\end{document}